%\documentclass{book}\usepackage{sjtfonts}\usepackage{includegraphics}\usepackage{alltt}\usepackage{latexsym}
%\usepackage{array}\usepackage{multicol}
%\begin{document}\input{defs0}%		miradefs.tex 		    June 1994

% opening and closing a quoted Miranda section

\newcommand{\mira}{\noindent\hrulefill\nopagebreak\so}
\newcommand{\arim}{\nopagebreak\st\hrulefill\\}

% backslash in the Miranda environment
% Only feasible approach is to use \verb+\+
% in line!

%\newcommand{\bs}{$\mi\backslash\!\!$}
%\newcommand{\lbs}{\mbox{\verb+\+}}

% ignoring part of the input

\newcommand{\ignore}[1]{}

% a blank line in a Miranda section

\newcommand{\bl}{\mbox{\ }}

% Signalling a diffculty.

\definecolor{light-gray}{gray}{0.95}

\newcommand{\beware}[2]
 {\medskip\noindent\colorbox{light-gray}{\parbox{\textwidth}
 {\mbox{\large\textbf{#1}} \medskip\noindent\\ #2}\medskip\noindent}}

% Subscripting in Miranda text
\newcommand{\subscr}[2]{\mbox{${\mbox{\texttt{#1}}}_{\mbox{\texttt{#2}}}$}}

%\newcommand{\subscr}[2]{\mbox{${\mbox{\mi #1}}_{\mbox{\smtt #2}}$}}
\newcommand{\aone}{\subscr{a}{1}}
\newcommand{\atwo}{\subscr{a}{2}}
\newcommand{\ak}{\subscr{a}{k}}

\newcommand{\aoneone}{\subscr{a}{1,1}}

\newcommand{\eone}{\subscr{e}{1}}
\newcommand{\etwo}{\subscr{e}{2}}
\newcommand{\ei}{\subscr{e}{i}}
\newcommand{\er}{\subscr{e}{r}}

\newcommand{\fone}{\subscr{f}{1}}
\newcommand{\fl}{\subscr{f}{l}}

\newcommand{\gone}{\subscr{g}{1}}
\newcommand{\gtwo}{\subscr{g}{2}}
\newcommand{\gi}{\subscr{g}{i}}
\newcommand{\gr}{\subscr{g}{r}}

\newcommand{\hone}{\subscr{h}{1}}
\newcommand{\hl}{\subscr{h}{l}}

\newcommand{\pone}{\subscr{p}{1}}
\newcommand{\ptwo}{\subscr{p}{2}}
\newcommand{\pk}{\subscr{p}{k}}

\newcommand{\qone}{\subscr{q}{1}}
\newcommand{\qtwo}{\subscr{q}{2}}
\newcommand{\qk}{\subscr{q}{k}}

\newcommand{\rone}{\subscr{r}{1}}
\newcommand{\rtwo}{\subscr{r}{2}}

\newcommand{\tone}{\subscr{t}{1}}
\newcommand{\ttwo}{\subscr{t}{2}}
\newcommand{\tn}{\subscr{t}{n}}

\newcommand{\vone}{\subscr{v}{1}}
\newcommand{\vtwo}{\subscr{v}{2}}
\newcommand{\vn}{\subscr{v}{n}}

% for regular expressions

\newcommand{\stone}{\subscr{st}{1}}
\newcommand{\sttwo}{\subscr{st}{2}}
\newcommand{\stn}{\subscr{st}{n}}
\newcommand{\rthree}{\subscr{r}{3}}

\newcommand{\eps}{$\varepsilon$}
\newcommand{\trans}[1]{\stackrel{\mbox{\tt #1}}{\longrightarrow}}



% superscript in Miranda text

\newcommand{\superscr}[2]{\mbox{${\mbox{\mi #1}}^{\mbox{\smtt #2}}$}}

% Not equals in Miranda text

\newcommand{\noteq}{\makebox[0pt][l]{=}/}
\newcommand{\overslash}{\makebox[0pt][l]{/}}

% Definitions in complexity chapter

\newfont{\oldSym}{cmsy10}
\newcommand{\Th}{\boldmath$\oldSym\Theta$}
\newcommand{\pow}[1]{\^{}#1}
\newcommand{\leqq}{$\leq$}
\newcommand{\geqq}{$\geq$}
\newcommand{\lll}{$\ll$}
\newcommand{\eqq}{$\equiv$}
\newcommand{\ltwo}{\subscr{log}{2}}

% Indexing

\newcommand{\minx}[1]{\index{#1@{\ttfamily{}#1}}}
\appendix


%\minitocoptionfalse
\chapter{Haskell operators}
\label{hsOps}
\index{operator!table of properties}
\index{floating point operators}
\normalfont\normalsize
The operators in the Haskell prelude are listed below in decreasing order of binding power:
see Section \ref{operators} for a discussion of associativity and binding power.

\begin{center}
%\psframebox[cornersize=absolute,linearc=6pt,linecolor=black,linewidth=0.5pt,framesep=0pt]
{\begin{tabular}{p{.3in}p{1in}p{1.5in}p{1in}}
\\[-9pt]
 & Left associative & Non-associative & Right associative\\
\hline
9 & \ttfamily !! & & . \\
8 & & & \ttfamily **, \^{}, \^{}\^{}\\
7 & \ttfamily   *, /, `div`, `mod`, `rem`, `quot`& \\
6 & \ttfamily +, - & & \ttfamily \\
5 & \ttfamily &  & \ttfamily :, ++ \\
4 & & \ttfamily /=, <, <=, ==, >, >=, `elem`, `notElem` \\
3 & & & \ttfamily \&\& \\
2 & & & \ttfamily || \\
1 & \ttfamily >>, >>= \\
0 & & & \ttfamily \$, \$!, `seq` \\[3pt]
\end{tabular}}
\end{center}
\minx{"!}\minx{"!"!}\minx{.}
\minx{**}\minx{\^{}}\minx{\^{}\^{}}
\minx{*}\minx{/}\index{mod@\texttt{`mod`}}\index{div@\texttt{`div`}}\index{rem@\texttt{`rem `}}\index{quot@\texttt{`quot `}}
\minx{+}\minx{-}
\minx{:}\minx{++}
\minx{/=}\minx{<}\minx{<=}\minx{==}\minx{>}\minx{>=}\index{elem@\texttt{`elem`}}\index{notElem@\texttt{`notElem`}}
\minx{\&\&}\minx{"|"|}\minx{>>}\minx{\$}\minx{\$"!}\index{seq@\texttt{`seq`}}
\medskip
\noindent
Also defined in this text are the operators
\index{\symbol{62}\symbol{62}\symbol{61}@\texttt{\symbol{62}\symbol{62}\symbol{61}}}
\index{\symbol{62}\symbol{46}\symbol{62}@\texttt{\symbol{62}\symbol{46}\symbol{62}}}
\index{\symbol{62}\symbol{42}\symbol{62}@\texttt{\symbol{62}\symbol{42}\symbol{62}}}
\begin{center}
%\psframebox[cornersize=absolute,linearc=6pt,linecolor=black,linewidth=0.5pt,framesep=0pt]
{\begin{tabular}{p{.3in}p{1in}p{1.5in}p{1in}}
\\[-9pt]
9 & \ttfamily  >.> \\ 
5 & & & \ttfamily >*> \\[3pt]
\end{tabular}}
\end{center}

\medskip
\noindent
The restrictions on names of operators, which are formed using the characters
\begin{alltt}
! \# \$ \% \& * + . / < = > ? \@ \bs{} \^{} | : - \twid
\end{alltt}
are that operators must not start with a colon; this character starts an infix
\emph{constructor}. The operators \texttt{-} and \texttt{!} can be user-defined, but note that they have
a special meaning in certain circumstances -- the obvious advice here is not to use them. 
Finally, certain combinations of symbols are reserved,
and cannot be used: {\ttfamily ..\ :\ ::\ => = @ \bs{} | \^{} <- ->}.



To change the associativity or binding power of an operator, {\ttfamily \&\&\&} say,
we make a declaration like\index{fixity!declaration}
\begin{alltt}
infixl 7 \&\&\&
\end{alltt}
which states that {\ttfamily \&\&\&} has binding power 7, and is a left associative
operator. We can also declare operators as non-associative ({\ttfamily infix}) and
right associative ({\ttfamily infixr}). Omitting the binding power gives a default
of 9. These declarations can also be used for back-quoted function names, as
in
\begin{alltt}
infix 0 `poodle`
\end{alltt}

%\end{document}
